\usepackage{makecell}
\usepackage{enumitem}
\usepackage{multirow}
\usepackage[para,online,flushleft]{threeparttable}
\usepackage{caption}
\geometry{left=30mm, right=15mm, top=20mm, bottom=20mm}


% Заголовки \chapter* по центру печатной области (в bmstu дают лишний \parindent слева)
\titleformat{name=\chapter,numberless}[block]{}{}{0pt}{\large\bfseries\centering}
\titlespacing{name=\chapter,numberless}{0pt}{-22pt}{10pt}

\DeclareNameAlias{byauthor}{family-given}
\ExecuteBibliographyOptions{
    maxbibnames=3,   % Максимум авторов до сокращения в библиографии
    minbibnames=3,   % Сокращаем до 3 имён + "и др."
}

% Работа с изображениями и таблицами; переопределение названий по ГОСТу
\captionsetup[table]{singlelinecheck=false, labelsep=endash}
\captionsetup[table]{justification=raggedright,singlelinecheck=off}
\captionsetup{labelsep=endash}
\captionsetup[figure]{name={Рисунок}}
\captionsetup[lstlisting]{justification=raggedright,singlelinecheck=off} % Работа с листингом

% Отточия у глав
\makeatletter
\renewcommand*{\l@chapter}[2]{
\ifnum \c@tocdepth>\m@ne
    \addpenalty{-\@highpenalty}
    \vskip 1em \@plus.2em
    \@dottedtocline{0}{0pt}{1.5em}{\bfseries #1}{\bfseries #2}
\fi
}
\makeatother

\renewcommand*{\arraystretch}{1} % математические матрицы интервал между строками

% \usepackage{ulem} % Нормальное нижнее подчеркивание
% \usepackage{hhline} % Двойная горизонтальная линия в таблицах
% \usepackage[figure,table]{totalcount} % Подсчет изображений, таблиц
% \usepackage{rotating} % Поворот изображения вместе с названием
% \usepackage{lastpage} % Для подсчета числа страниц

% Дополнительное окружения для подписей
% \usepackage{array}
% \newenvironment{signstabular}[1][1]{
%       \renewcommand*{\arraystretch}{#1}
%       \tabular
% }{
%       \endtabular
% }

% \usepackage{pgfplots}
% \usetikzlibrary{datavisualization}
% \usetikzlibrary{datavisualization.formats.functions}

\usepackage[justification=centering]{caption} % Настройка подписей float объектов

\newcommand{\code}[1]{\texttt{#1}}

\usepackage{listings}
% Для листинга кода:
\lstset{%
        language=python,                                           % выбор языка для подсветки
        basicstyle=\small\sffamily,                     % размер и начертание шрифта для подсветки кода
        numbersep=5pt,
        numbers=none,                                           % где поставить нумерацию строк (слева\справа)
        %numberstyle=,                                      % размер шрифта для номеров строк
        stepnumber=1,                                           % размер шага между двумя номерами строк
        xleftmargin=0pt,
        showstringspaces=false,
        numbersep=5pt,                                          % как далеко отстоят номера строк от подсвечиваемого кода
        frame=single,                                           % рисовать рамку вокруг кода
        tabsize=4,                                                      % размер табуляции по умолчанию равен 4 пробелам
        captionpos=t,                                           % позиция заголовка вверху [t] или внизу [b]
        breaklines=true,
        breakatwhitespace=true,                         % переносить строки только если есть пробел
        escapeinside={\#*}{*)},                         % если нужно добавить комментарии в коде
        backgroundcolor=\color{white}
}


\usepackage{pdfpages}
\usepackage{totcount}
\usepackage{hyperref}
\usepackage{lastpage}

\lstset{
        literate=
        {а}{{\selectfont\char224}}1
        {б}{{\selectfont\char225}}1
        {в}{{\selectfont\char226}}1
        {г}{{\selectfont\char227}}1
        {д}{{\selectfont\char228}}1
        {е}{{\selectfont\char229}}1
        {ё}{{\"e}}1
        {ж}{{\selectfont\char230}}1
        {з}{{\selectfont\char231}}1
        {и}{{\selectfont\char232}}1
        {й}{{\selectfont\char233}}1
        {к}{{\selectfont\char234}}1
        {л}{{\selectfont\char235}}1
        {м}{{\selectfont\char236}}1
        {н}{{\selectfont\char237}}1
        {о}{{\selectfont\char238}}1
        {п}{{\selectfont\char239}}1
        {р}{{\selectfont\char240}}1
        {с}{{\selectfont\char241}}1
        {т}{{\selectfont\char242}}1
        {у}{{\selectfont\char243}}1
        {ф}{{\selectfont\char244}}1
        {х}{{\selectfont\char245}}1
        {ц}{{\selectfont\char246}}1
        {ч}{{\selectfont\char247}}1
        {ш}{{\selectfont\char248}}1
        {щ}{{\selectfont\char249}}1
        {ъ}{{\selectfont\char250}}1
        {ы}{{\selectfont\char251}}1
        {ь}{{\selectfont\char252}}1
        {э}{{\selectfont\char253}}1
        {ю}{{\selectfont\char254}}1
        {я}{{\selectfont\char255}}1
        {А}{{\selectfont\char192}}1
        {Б}{{\selectfont\char193}}1
        {В}{{\selectfont\char194}}1
        {Г}{{\selectfont\char195}}1
        {Д}{{\selectfont\char196}}1
        {Е}{{\selectfont\char197}}1
        {Ё}{{\"E}}1
        {Ж}{{\selectfont\char198}}1
        {З}{{\selectfont\char199}}1
        {И}{{\selectfont\char200}}1
        {Й}{{\selectfont\char201}}1
        {К}{{\selectfont\char202}}1
        {Л}{{\selectfont\char203}}1
        {М}{{\selectfont\char204}}1
        {Н}{{\selectfont\char205}}1
        {О}{{\selectfont\char206}}1
        {П}{{\selectfont\char207}}1
        {Р}{{\selectfont\char208}}1
        {С}{{\selectfont\char209}}1
        {Т}{{\selectfont\char210}}1
        {У}{{\selectfont\char211}}1
        {Ф}{{\selectfont\char212}}1
        {Х}{{\selectfont\char213}}1
        {Ц}{{\selectfont\char214}}1
        {Ч}{{\selectfont\char215}}1
        {Ш}{{\selectfont\char216}}1
        {Щ}{{\selectfont\char217}}1
        {Ъ}{{\selectfont\char218}}1
        {Ы}{{\selectfont\char219}}1
        {Ь}{{\selectfont\char220}}1
        {Э}{{\selectfont\char221}}1
        {Ю}{{\selectfont\char222}}1
        {Я}{{\selectfont\char223}}1
}

\usepackage[listings]{tcolorbox}
\tcbuselibrary{minted, breakable, skins}

\newtcbinputlisting[auto counter,number within=chapter]{\mylisting}[4][]{
        enhanced,
        before skip=10pt,
        arc=0mm,
        boxrule=1pt,
        listing only,
        listing file={#1},
        listing engine=listings,
        listing options={
                language=python,                                           % выбор языка для подсветки
                basicstyle=\small\sffamily,                     % размер и начертание шрифта для подсветки кода
                numbersep=5pt,
                numbers=left,
                numberstyle=\tiny,
                breaklines=true,
                numbersep=1.5mm,
                % numbers=none,                                           % где поставить нумерацию строк (слева\справа)
                %numberstyle=,                                      % размер шрифта для номеров строк
                stepnumber=1,                                           % размер шага между двумя номерами строк
                xleftmargin=0pt,
                showstringspaces=false,
                numbersep=5pt,                                          % как далеко отстоят номера строк от подсвечиваемого кода
                frame=none,                                           % рисовать рамку вокруг кода
                tabsize=4,                                                      % размер табуляции по умолчанию равен 4 пробелам
                captionpos=t,                                           % позиция заголовка вверху [t] или внизу [b]
                breaklines=true,
                breakatwhitespace=true,                         % переносить строки только если есть пробел
                escapeinside={\#*}{*)},                         % если нужно добавить комментарии в коде
                backgroundcolor=\color{white}
                #3
        },
        colback=white,
        colframe=black,
        breakable,
        title={\parbox{\linewidth}{\raggedright Листинг~\thetcbcounter~--~#2}},
        attach boxed title to top left,
        boxed title style={
                enhanced jigsaw,
                opacityfill=1,
                sharp corners,
                boxrule=0pt,
                left=-3pt,
                right=0pt,
                width=\linewidth,                  % разрешаем перенос заголовка
        },
        coltitle=black,
        colbacktitle=white,
        enlarge top at break by=5mm,
        overlay first={
                \draw[line width=.5pt] (frame.south west)--(frame.south east);
        },
        overlay middle={
                \draw[line width=.5pt] (frame.south west)--(frame.south east);
                \draw[line width=.5pt] (frame.north west)--(frame.north east);
                \node[anchor=south west, inner xsep=0pt, xshift=0pt] at (frame.north west) {Продолжение листинга \thetcbcounter};
        },
        overlay last={
                \draw[line width=.5pt] (frame.north west)--(frame.north east);
                \node[anchor=south west, inner xsep=0pt, xshift=0pt] at (frame.north west) {Продолжение листинга \thetcbcounter};
        },
        #4
}
